\documentclass[12pt,a4paper]{article}

\usepackage[utf8]{inputenc}
\usepackage[T1]{fontenc}
\usepackage{lmodern}
\usepackage{amsmath,amssymb,amsthm,mathtools}
\usepackage{geometry}
\usepackage{enumitem}
\usepackage{hyperref}
\usepackage{cleveref}
\usepackage{booktabs}
\usepackage{xcolor}
\usepackage{fancyhdr}

\geometry{margin=2.5cm}

\newtheoremstyle{postulate}
  {12pt}{12pt}{\itshape}{0pt}{\bfseries}{.}{.5em}{}
\theoremstyle{postulate}
\newtheorem{postulate}{Postulate}
\newtheorem{definition}{Definition}[section]
\newtheorem{remark}{Remark}[section]
\newtheorem{proposition}{Proposition}[section]

\newcommand{\Tr}{\operatorname{Tr}}
\newcommand{\spec}{\operatorname{spec}}
\newcommand{\Hil}{\mathcal{H}}
\newcommand{\Alg}{\mathcal{A}}
\newcommand{\Dirac}{D}
\newcommand{\Cut}{\Lambda}
\newcommand{\dd}{\mathrm{d}}
\newcommand{\Causal}{\mathcal{C}}
\newcommand{\Kset}{\mathfrak{K}}
\newcommand{\planck}{\ell_{\mathrm P}}
\newcommand{\zetaD}{\zeta_D}

\pagestyle{fancy}
\fancyhf{}
\fancyhead[L]{\small\textsc{SCT Theory --- Axiom Probe}}
\fancyhead[R]{\small\thepage}
\renewcommand{\headrulewidth}{0.4pt}

\hypersetup{
  colorlinks=true,
  linkcolor=blue!60!black,
  citecolor=green!40!black,
  urlcolor=blue!60!black,
  pdftitle={Decorated Causal Order as Kinematics in SCT},
  pdfauthor={David Alfyorov},
}

\title{\textbf{Decorated Causal Order as Kinematics in SCT}\\[6pt]
\large An Axiom-Style Probe on Spectral Data over a Locally Finite Causal Set}
\author{David Alfyorov}
\date{March 2026\\[4pt]\small Working Note}

\begin{document}
\maketitle

\begin{center}
\small\textit{Credits: Aliaksandr Samatyia contributed research-assistance and workflow support.}
\end{center}

\thispagestyle{empty}

\begin{abstract}
This note is written as an axiom-style probe in the house style of the SCT
foundational corpus.  The aim is not to enlarge the canonical postulate list,
but to test one kinematic statement in the same register: a bare causal set
is too poor to support the full geometric and internal content expected of a
quantum-gravitational theory, while a smooth manifold is too rich to be
assumed at the fundamental level.  A natural intermediate kinematic object is
a locally finite causal order whose events carry spectral data.

We formulate this idea as a single postulate, define the required
mathematical objects, explain the physical motivation, and describe the
limiting regimes in which ordinary spacetime geometry may emerge.  The text
is intentionally conservative.  It does not claim that the full microscopic
dynamics is already known.  It isolates only the kinematic content that would
make such a dynamics meaningful.
\end{abstract}

\tableofcontents
\newpage

\section{Preliminary Definitions}
\label{sec:prelim}

Before stating the postulate, we fix the objects that will appear in it.

\begin{definition}[Locally finite causal set]
A causal set is a pair $(\Causal,\preceq)$ where $\Causal$ is a set and
$\preceq$ is a partial order such that:
\begin{enumerate}[label=(\roman*)]
  \item $x \preceq x$ for all $x \in \Causal$;
  \item $x \preceq y$ and $y \preceq x$ imply $x=y$;
  \item $x \preceq y$ and $y \preceq z$ imply $x \preceq z$;
  \item for any $x,z\in\Causal$, the interval
  \[
    I(x,z)=\{y\in\Causal\mid x\preceq y\preceq z\}
  \]
  is finite.
\end{enumerate}
\end{definition}

\begin{definition}[Spectral triple]
A spectral triple is a datum $(\Alg,\Hil,\Dirac)$ consisting of an algebra
$\Alg$ represented on a Hilbert space $\Hil$ together with a self-adjoint
operator $\Dirac$ with compact resolvent and bounded commutators
$[\Dirac,a]$ for all $a\in\Alg$.
\end{definition}

\begin{definition}[Spectral zeta function]
Given a Dirac operator $\Dirac$ with nonzero eigenvalues $\{\lambda_n\}$, its
spectral zeta function is
\begin{equation}
  \zetaD(s) = \sum_{\lambda_n\neq 0} |\lambda_n|^{-s},
  \label{eq:zeta}
\end{equation}
for $\operatorname{Re}(s)$ large enough, with meromorphic continuation
elsewhere when such continuation exists.
\end{definition}

\begin{definition}[Connes distance]
For a commutative spectral triple, the spectral distance is
\begin{equation}
  d_{\mathrm C}(p,q)
  = \sup\left\{
      |a(p)-a(q)|
      \;\middle|\;
      a\in\Alg,\ \|[\Dirac,a]\|\leq 1
    \right\}.
  \label{eq:connes-distance}
\end{equation}
\end{definition}

\begin{definition}[Spectral dimension]
If the heat trace of $\Dirac^2$ is available, one defines the scale-dependent
spectral dimension by
\begin{equation}
  d_S(\sigma)
  = -2\,\frac{\dd \ln \Tr(e^{-\sigma \Dirac^2})}{\dd \ln \sigma}.
  \label{eq:spectral-dimension}
\end{equation}
\end{definition}

\begin{definition}[Decorated event]
A decorated event is an element $x\in\Causal$ endowed with local spectral
data
\begin{equation}
  \mathfrak{s}(x) = (\Alg_x,\Hil_x,\Dirac_x).
  \label{eq:decorated-event}
\end{equation}
\end{definition}

\section{Postulate 1 --- Decorated Causal Kinematics}
\label{sec:postulate}

\begin{postulate}[Decorated causal kinematics]
The fundamental kinematic data of SCT is a quadruple
\begin{equation}
  \Kset
  = \left(
      \Causal,\preceq,
      \{\mathfrak{s}(x)\}_{x\in\Causal},
      \{\Phi_{x\to y}\}_{x\prec y}
    \right),
  \label{eq:kset}
\end{equation}
where:
\begin{enumerate}[label=(\roman*)]
  \item $(\Causal,\preceq)$ is a locally finite causal set;
  \item each event $x\in\Causal$ carries a spectral triple
  $\mathfrak{s}(x)=(\Alg_x,\Hil_x,\Dirac_x)$;
  \item each causal relation $x\prec y$ carries a comparison map
  $\Phi_{x\to y}:\Hil_x\to\Hil_y$;
  \item smooth spacetime geometry is not fundamental but emergent from the
  collective spectral-causal data in an appropriate continuum limit.
\end{enumerate}
\end{postulate}

\paragraph{Immediate reading.}
The postulate says that neither causal order alone nor spectral geometry
alone is taken as the whole microscopic story.  Order captures precedence and
local finiteness; spectral data carries the local geometric, matter, and
operator-theoretic content from which metric notions may later emerge.

\paragraph{Kinematic rather than dynamical status.}
No dynamical law has yet been stated.  The postulate concerns the space of
microscopic configurations, not their weighting.  That distinction is
essential.  In the logic of the theory, one must specify the kinematic arena
before assigning a spectral action or renormalization-group flow on it.

\paragraph{Immediate formal consequence.}
If an effective commutative limit exists, the postulate already singles out
the correct kind of data from which both causal and spectral notions of
distance might be reconstructed.  It therefore has mathematical substance
even before a microscopic action is written down.

\section{Mathematical Content}
\label{sec:math-content}

\subsection{Why causal order alone is insufficient}

A pure causal set furnishes chronology and local finiteness, but it does not
by itself specify:
\begin{enumerate}[label=(\alph*)]
  \item internal gauge structure,
  \item fermionic data,
  \item spectral dimension,
  \item the operatorial information required for a spectral action.
\end{enumerate}
Thus a causal set by itself is an ordering structure, not yet a geometric or
field-theoretic object in the sense required by SCT.

\subsection{Why a smooth manifold is too strong}

At the opposite extreme, a differentiable manifold with metric already
assumes the existence of the continuum geometry that quantum gravity is meant
to explain.  Near the Planck scale, where one expects discreteness,
nonlocality, or a change of effective dimension, starting from a smooth
metric manifold builds too much infrared structure into the ultraviolet
foundation.

\subsection{Role of the comparison maps}

The maps $\Phi_{x\to y}$ are not introduced as a finished microscopic
propagator.  They are the minimal kinematic glue needed to compare spectral
data across causally related events.  At the very least, one expects some
compatibility with causal composition,
\begin{equation}
  \Phi_{x\to z}
  \sim \Phi_{y\to z}\circ \Phi_{x\to y}
  \qquad\text{whenever } x\prec y\prec z,
  \label{eq:composition}
\end{equation}
though the precise relation may be modified by coarse-graining, stochastic
structure, or quantum interference in the full microscopic theory.

\subsection{Spectral reconstruction target}

The continuum target of the postulate is not arbitrary.  In a commutative
limit one expects the effective algebra to approach
\begin{equation}
  \Alg_{\mathrm{eff}} \simeq C^\infty(M),
  \label{eq:commutative}
\end{equation}
so that ordinary geometric quantities are reconstructed spectrally from an
effective Dirac operator.  The postulate therefore names not only the
microscopic arena but also the shape of the desired infrared recovery.

\section{What the Postulate Does and Does Not Claim}
\label{sec:nonclaims}

It is important to state explicitly what is and is not being asserted.

\begin{enumerate}[label=(\alph*)]
  \item The postulate \emph{does} propose a candidate microscopic kinematic
  arena.
  \item It \emph{does not} yet provide the weighting of microscopic
  configurations.
  \item It \emph{does not} provide a rigorous continuum-limit theorem.
  \item It \emph{does} require that any later dynamics act on or be
  equivalent to some decorated causal-spectral structure of this kind.
\end{enumerate}

The point of such explicit discipline is methodological.  Foundational texts
in quantum gravity often become ambiguous precisely where kinematics and
dynamics are allowed to blur.  The SCT methodology requires keeping them
separate.

\section{Physical Motivation}
\label{sec:motivation}

\subsection{General relativity points toward causal primacy}

Lorentzian geometry teaches that causal structure is not a minor decoration
of spacetime.  In many settings, the conformal class of the metric is already
encoded in the causal order.  This suggests that any ultraviolet completion
which ignores causal structure at the foundational level is discarding
something too important.

\subsection{Noncommutative geometry points toward spectral primacy}

Connes' noncommutative geometry teaches a complementary lesson: geometry can
be encoded in operator spectra rather than in coordinate charts.  Distances,
dimensions, and curvature information are read from $\Dirac$ and the algebra
acting on its Hilbert space.

\subsection{SCT reads both lessons together}

The kinematic wager of SCT is therefore not to choose between causal order
and spectral data, but to hold them together from the start.  Causal order
without spectral content is too poor.  Spectral content without causal order
risks losing the Lorentzian backbone required for spacetime physics.

\paragraph{Discipline about status.}
This is not yet a theorem.  It is a foundational proposal motivated by
structural pressure from both sides.  Its seriousness depends on whether it
makes the later dynamical and phenomenological questions sharper rather than
vaguer.

\section{Formal Consequences if the Postulate is Accepted}
\label{sec:consequences}

Accepting the postulate commits the later theory to several downstream
requirements.

\begin{proposition}[Continuum compatibility requirement]
Any acceptable dynamics on $\Kset$ must admit an infrared regime in which the
effective spectral data defines distances, dimensions, and curvature notions
compatible with known low-energy geometry.
\end{proposition}

\begin{proof}
If no such regime exists, the postulate fails the correspondence test with
established gravitational physics.  Since SCT is intended as a candidate UV
completion rather than a rejection of the infrared world, such a regime is
mandatory.
\end{proof}

\begin{proposition}[Observability requirement]
Any physically meaningful realization of the postulate must ultimately define
observables built from either spectral invariants, causal-order invariants,
or mixed quantities derived from both.
\end{proposition}

\begin{proof}
A kinematic construction with no observable sector would not qualify as a
physical theory candidate.  The postulate therefore imposes an observability
burden on every later dynamical completion.
\end{proof}

\section{Limiting Cases}
\label{sec:limits}

Any foundational postulate should be judged by the limiting regimes it makes
plausible.  Four are especially relevant.

\subsection{Classical geometric limit}

In a suitable coarse-grained regime, one expects an effective commutative
spectral triple
\begin{equation}
  (\Alg_{\mathrm{eff}},\Hil_{\mathrm{eff}},\Dirac_{\mathrm{eff}})
  \simeq
  \bigl(C^\infty(M),L^2(S),\Dirac_M\bigr),
  \label{eq:classical-limit}
\end{equation}
from which ordinary Riemannian geometry can be reconstructed.  The postulate
is acceptable only if such a recovery channel exists.

\subsection{Weak-field, low-curvature regime}

At scales far above the microscopic discreteness length and in a regime of
weak curvature, the spectral action built from the emergent Dirac operator
must reduce to the Einstein--Hilbert term plus controlled higher-curvature
corrections:
\begin{equation}
  \Tr f\!\left(\frac{\Dirac_{\mathrm{eff}}^2}{\Cut^2}\right)
  \Longrightarrow
  \int \sqrt{g}\,
  \left(
    \frac{R-2\Lambda}{16\pi G}
    + \alpha R^2 + \beta R_{\mu\nu}R^{\mu\nu}
    + \cdots
  \right)\dd^4x.
  \label{eq:weak-field-limit}
\end{equation}
If the postulate cannot support this passage, it does not connect to known
gravitational physics.

\subsection{Large-distance regime}

At sufficiently large scales one expects the microscopic discreteness to be
washed out.  The causal order should approximate the causal structure of a
smooth Lorentzian manifold, while the local spectral data should glue
together into the familiar infrared geometry and matter content.

\subsection{Planckian regime}

Near the Planck scale, one expects the opposite tendency.  The discrete
nature of $\Causal$ and the genuinely spectral nature of the geometry should
become visible.  In that regime it is natural to expect scale-dependent
spectral dimension, breakdown of naive manifold locality, and a more
primitive operatorial notion of geometry than the metric tensor.

\subsection{Interpretive consequence of the four limits}

Taken together, the four limits say something restrictive.  The postulate is
not free to remain a purely abstract combinatorial construction.  It must
interpolate between a genuinely non-manifold ultraviolet regime and a
recognizable low-energy geometric regime.  That is exactly the type of
tension any viable quantum-gravitational kinematics should be able to
survive.

\section{Constructive Hints from Experiment and Established Physics}
\label{sec:hints}

The postulate is motivated not only by mathematical aesthetics but by a set
of convergent hints from known physics.

\begin{table}[ht]
\centering
\renewcommand{\arraystretch}{1.25}
\begin{tabular}{@{}p{0.40\textwidth}p{0.10\textwidth}p{0.40\textwidth}@{}}
\toprule
\textbf{Hint} & \textbf{Use} & \textbf{Structural push} \\
\midrule
Causal structure already controls large parts of Lorentzian geometry & causal & take order seriously at the kinematic level \\
\addlinespace
Connes reconstruction and spectral geometry & spectral & encode geometry through operator data rather than coordinates \\
\addlinespace
Area/entropy relations in gravity & mixed & support non-metric geometric primitives at fundamental level \\
\addlinespace
Dimensional flow in quantum-gravity approaches & UV & make room for a non-manifold microscopic regime \\
\addlinespace
Entanglement/geometry connections in holography and QFT & relational & suggest that geometric proximity should emerge from deeper quantum structure \\
\bottomrule
\end{tabular}
\caption{Constructive hints that make decorated causal kinematics a serious
candidate starting point rather than a purely decorative proposal.}
\label{tab:hints}
\end{table}

These hints do not amount to a proof of the postulate.  Their role is
different: they show that the postulate is not arbitrary.  It sits at a
crossing point of several independent lessons from established physics and
quantum-gravity research.

\section{Relation to the Rest of the Theory}
\label{sec:relation}

This postulate is upstream of several later tasks.
\begin{enumerate}[label=(\roman*)]
  \item One still needs a dynamical principle, for instance a spectral action
  or equivalent weighting functional on the decorated causal data.
  \item One still needs a coarse-graining map that identifies the emergent
  continuum observables.
  \item One still needs a precise relation between the comparison maps
  $\Phi_{x\to y}$ and any microscopic entanglement or transport structure.
\end{enumerate}

That incompleteness is not a flaw of the postulate itself.  It is simply the
normal boundary between kinematics and dynamics.  What matters is that the
kinematic proposal makes those later questions meaningful rather than
ill-posed.

\paragraph{Relation to spectral action dynamics.}
If one later assigns an effective spectral action of the form
\begin{equation}
  S_{\mathrm{sp}} = \Tr f\!\left(\frac{\Dirac^2}{\Cut^2}\right),
  \label{eq:spectral-action}
\end{equation}
then the present postulate supplies the natural microscopic arena on which
such a dynamics may be interpreted after coarse graining.

\section{Consistency and Falsifiability Burden}
\label{sec:burden}

Any foundational postulate must carry a burden, not only an ambition.  For
the present proposal, at least four consistency requirements are unavoidable.

\begin{enumerate}[label=(\roman*)]
  \item \textbf{Lorentz-compatibility burden.} The causal order must admit an
  infrared regime compatible with Lorentzian spacetime rather than secretly
  selecting a rigid preferred lattice frame.
  \item \textbf{Spectral-reconstruction burden.} The decorated data must be
  rich enough to recover effective distances, dimension, and curvature in a
  controlled continuum limit.
  \item \textbf{Dynamics burden.} There must exist a natural action or
  transition rule on $\Kset$, or the postulate remains a purely static
  ontology.
  \item \textbf{Observability burden.} The theory must identify observables
  that depend on the spectral-causal data and can, at least in principle, be
  confronted with experiment or observation.
\end{enumerate}

These burdens are not objections.  They are the proper scientific price of a
serious foundational proposal.

\paragraph{What could falsify the kinematic picture.}
The postulate would be in grave difficulty if it turned out that:
\begin{enumerate}[label=(\alph*)]
  \item no continuum regime with acceptable infrared geometry exists,
  \item the spectral-causal data cannot reproduce known low-energy matter and
  gravitational structure,
  \item the required observables remain permanently undefined,
  \item the microscopic construction necessarily violates empirical
  Lorentz-invariance bounds.
\end{enumerate}
In that sense the postulate is not protected by vagueness; it remains open to
failure on concrete physical grounds.

\section{Summary and Open Problems}
\label{sec:summary}

\subsection{Summary}

The content of the note can be stated in one sentence: the fundamental
kinematic arena of SCT may be taken to be a locally finite causal order whose
events carry spectral triples and are connected by comparison maps.

This proposal tries to preserve, at once, two lessons:
\begin{enumerate}[label=(\roman*)]
  \item causal order is too fundamental to be demoted to a derived afterthought;
  \item spectral data is rich enough to encode geometry and matter.
\end{enumerate}
The point of the postulate is therefore not ornamental complexity.  It is the
attempt to avoid choosing too little structure on one side and too much
structure on the other.

\subsection{Postulate summary table}

\begin{center}
\renewcommand{\arraystretch}{1.2}
\begin{tabular}{@{}p{3.0cm}p{10.0cm}@{}}
\toprule
\textbf{Component} & \textbf{Content} \\
\midrule
Causal substrate & Locally finite partially ordered set $(\Causal,\preceq)$ \\
Local geometric data & Spectral triple $\mathfrak{s}(x)=(\Alg_x,\Hil_x,\Dirac_x)$ at each event \\
Inter-event structure & Comparison maps $\Phi_{x\to y}$ along causal relations \\
Continuum target & Effective commutative spectral geometry in the infrared \\
Scientific burden & Must support dynamics, observables, and low-energy recovery \\
\bottomrule
\end{tabular}
\end{center}

\subsection{Constructive program}

If one takes the postulate seriously, the next mathematical tasks are also
clear:
\begin{enumerate}[label=\arabic*.]
  \item define an admissible class of comparison maps $\Phi_{x\to y}$,
  \item specify a coarse-graining prescription from $\Kset$ to an effective
  spectral triple,
  \item identify invariant observables on the microscopic decorated causal
  data,
  \item formulate a dynamics that acts naturally on that kinematic arena.
\end{enumerate}
This constructive program is exactly what separates a promising foundational
ansatz from a merely evocative slogan.

\subsection{Open problems}

Several problems remain genuinely open:
\begin{enumerate}[label=\arabic*.]
  \item What is the best Lorentzian analogue of the spectral-triple language
  in the fully causal microscopic regime?
  \item Are the comparison maps $\Phi_{x\to y}$ fundamental objects, or only
  effective summaries of a deeper entanglement structure?
  \item Which universality classes of decorated causal sets flow to the same
  infrared geometry?
  \item Can one derive the continuum spectral action and its field content
  from a microscopic dynamics on $\Kset$ rather than imposing that relation
  from above?
  \item Which observables most directly falsify the decorated-causal
  kinematic picture?
\end{enumerate}

These are not marginal details.  They determine whether the postulate remains
an interesting kinematic ansatz or grows into a genuine candidate framework.

\begin{thebibliography}{99}

\bibitem{Bombelli1987}
L.~Bombelli, J.~Lee, D.~Meyer and R.~Sorkin,
``Space-time as a causal set,''
\emph{Phys. Rev. Lett.} \textbf{59} (1987) 521--524.

\bibitem{Connes1994Axiom}
A.~Connes,
\emph{Noncommutative Geometry},
Academic Press, 1994.

\bibitem{Connes2008Reconstruction}
A.~Connes,
``On the spectral characterization of manifolds,''
\emph{J. Noncommut. Geom.} \textbf{7} (2013) 1--82,
arXiv:\href{https://arxiv.org/abs/0810.2088}{0810.2088}.

\bibitem{Sorkin2003}
R.~D.~Sorkin,
``Causal sets: Discrete gravity,''
in \emph{Lectures on Quantum Gravity}, edited by A.~Gomberoff and
D.~Marolf, Springer, 2005,
arXiv:\href{https://arxiv.org/abs/gr-qc/0309009}{gr-qc/0309009}.

\bibitem{Vassilevich2003Axiom}
D.~V.~Vassilevich,
``Heat kernel expansion: User's manual,''
\emph{Phys. Rept.} \textbf{388} (2003) 279--360,
arXiv:\href{https://arxiv.org/abs/hep-th/0306138}{hep-th/0306138}.

\end{thebibliography}

\end{document}
